% CSE 715 — Computer Graphics
% Remaining real-yield topics: Point-obscures-point visibility, Back-face/visible-surface
% definitions, Trapezoid-primitive polygon fill + vertex/edge/surface matrices
% University of Chittagong, 7th Semester B.Sc. Engineering
%
% Compile with: pdflatex RemainingTopics_Solutions.tex (run twice for TOC)

\documentclass[12pt, a4paper]{article}

\usepackage[left=2.5cm, right=2.5cm, top=2.8cm, bottom=2.8cm, headheight=15pt]{geometry}
\usepackage{amsmath, amssymb, mathtools}
\usepackage{tikz}
\usetikzlibrary{arrows.meta, positioning, shapes.geometric, calc, backgrounds, fit}
\usepackage{booktabs}
\usepackage{array}
\usepackage{xcolor}
\usepackage{enumitem}
\usepackage{fancyhdr}
\usepackage{titlesec}
\usepackage{mdframed}
\usepackage{tabularx}
\usepackage{hyperref}

\definecolor{sectionblue}{RGB}{10,60,110}
\definecolor{boxbg}{RGB}{242,247,255}
\definecolor{answerbg}{RGB}{240,255,240}
\definecolor{notesbg}{RGB}{255,250,230}

\mdfdefinestyle{solutionframe}{linecolor=sectionblue, backgroundcolor=boxbg, roundcorner=4pt, linewidth=1.4pt, innertopmargin=7pt, innerbottommargin=7pt, innerleftmargin=9pt, innerrightmargin=9pt}
\mdfdefinestyle{answerframe}{linecolor=green!55!black, backgroundcolor=answerbg, roundcorner=4pt, linewidth=1.4pt, innertopmargin=7pt, innerbottommargin=7pt, innerleftmargin=9pt, innerrightmargin=9pt}
\mdfdefinestyle{notesframe}{linecolor=orange!70!black, backgroundcolor=notesbg, roundcorner=4pt, linewidth=1pt, innertopmargin=5pt, innerbottommargin=5pt, innerleftmargin=8pt, innerrightmargin=8pt}
\newenvironment{solutionbox}{\begin{mdframed}[style=solutionframe]}{\end{mdframed}}
\newenvironment{answerbox}{\begin{mdframed}[style=answerframe]}{\end{mdframed}}
\newenvironment{notesbox}{\begin{mdframed}[style=notesframe]}{\end{mdframed}}

\pagestyle{fancy}
\fancyhf{}
\lhead{\small\color{sectionblue} CSE 715 --- Remaining Topics}
\rhead{\small\color{sectionblue} Point-Obscures / Back-Face / Trapezoid Fill}
\cfoot{\small\thepage}
\renewcommand{\headrulewidth}{0.5pt}

\titleformat{\section}{\Large\bfseries\color{sectionblue}}{}{0em}{}[\vspace{-2pt}\color{sectionblue}\rule{\linewidth}{0.8pt}]
\titleformat{\subsection}{\large\bfseries\color{sectionblue!85!black}}{}{0em}{}

\begin{document}

\begin{center}
  {\LARGE\bfseries\color{sectionblue} CSE 715: Computer Graphics}\\[0.6em]
  {\Large\bfseries Point-Obscures-Point, Back-Face/Visible-Surface, Trapezoid-Primitive Fill}\\[0.3em]
  {\normalsize University of Chittagong \quad|\quad 7\textsuperscript{th} Semester B.Sc.\ (Engg.) Examination}
\end{center}

\vspace{0.3em}\noindent\rule{\linewidth}{1pt}\vspace{0.5em}

\tableofcontents
\newpage

%=====================================================================
\section{Point-Obscures-Point Visibility via Viewpoint (Ch10, Hidden Surfaces)}
%=====================================================================
\begin{notesbox}
\textbf{Re-verified against all 4 raw scans 2026-07-07} (this topic appears 2020, 2021, 2022, 2023 — not 2024). \textbf{Method:} compute $\mathbf{P}_i-\mathbf{V}$ for each point. Two points share a line of sight from $\mathbf{V}$ iff their $(\mathbf{P}-\mathbf{V})$ vectors are \textbf{parallel with the same sign} (one is a positive scalar multiple of the other) --- exactly the ray equation $\mathbf{r}(t)=\mathbf{V}+t\mathbf{d}$ from [[wiki/ray-tracing|Ch12]] evaluated at different $t$. Among collinear points, the \textbf{smaller} $t$ (nearer $\mathbf{V}$) obscures the larger. If no two points are collinear, none obscures any other --- all are independently visible. \textbf{2 of the 4 years actually have no occlusion at all} --- don't assume every year produces a clean "obscures" chain.
\end{notesbox}

\subsection*{2020 Q8(c) [3.5 marks]}
\begin{solutionbox}
Given points $P_1(1,2,0)$, $P_2(3,6,20)$, $P_3(2,4,10)$ and viewpoint $C(0,0,-10)$, determine which points obscure the others when viewed from $C$.
\end{solutionbox}
\begin{answerbox}
$P_1-C=(1,2,10)$, $P_2-C=(3,6,30)=3(1,2,10)$, $P_3-C=(2,4,20)=2(1,2,10)$ --- \textbf{all three are collinear} with $C$, all positive multiples of the same direction. Ordering by distance ($t=1,2,3$ respectively): $P_1$ is nearest, then $P_3$, then $P_2$. \textbf{$P_1$ obscures $P_3$, which in turn obscures $P_2$} --- all three lie on one ray, stacked front-to-back as $P_1,P_3,P_2$.
\end{answerbox}

\subsection*{2021 Q6(a) [3 marks]}
\begin{solutionbox}
Given points $P_1(1,3,1)$, $P_2(3,6,-11)$, $P_3(2,6,-5)$ and viewpoint $C(0,0,7)$, determine which points obscure the others when viewed from $C$.
\end{solutionbox}
\begin{answerbox}
$P_1-C=(1,3,-6)$, $P_2-C=(3,6,-18)$, $P_3-C=(2,6,-12)=2(1,3,-6)$. \textbf{$P_3$ is collinear with $P_1$} (exactly twice as far). Check $P_2$: $(3,6,-18)$ vs.\ $3(1,3,-6)=(3,9,-18)$ --- $y$-component doesn't match, so $P_2$ is \textbf{not} on the same line of sight as $P_1$/$P_3$.
\textbf{Conclusion: $P_1$ obscures $P_3$} (same ray, $P_1$ nearer). \textbf{$P_2$ is on a different line of sight} --- independently visible, neither obscuring nor obscured.
\end{answerbox}

\subsection*{2022 Q6(a) [3 marks]}
\begin{solutionbox}
Given points $A(1,5,-2)$, $B(1,4,9)$, $C(-2,-2,-3)$ and viewpoint $V(0,1,-10)$, determine which points obscure the others when viewed from $V$.
\end{solutionbox}
\begin{answerbox}
$A-V=(1,4,8)$, $B-V=(1,3,19)$, $C-V=(-2,-3,7)$. Checking all three pairs for a scalar-multiple relationship: none match (e.g.\ $A$ vs $B$: ratios $1, 1.33, 2.375$ --- inconsistent; $A$ vs $C$: ratios $-0.5,-1.33,1.14$ --- inconsistent; $B$ vs $C$: ratios $-0.5,-1,2.71$ --- inconsistent).
\textbf{Conclusion: no two of $A,B,C$ share a line of sight from $V$ --- none obscures any other. All three are simultaneously visible.}
\end{answerbox}

\subsection*{2023 Q6(a) [3 marks]}
\begin{solutionbox}
Given points $A(5,1,-2)$, $B(10,4,9)$, $C(15,-2,-3)$ and viewpoint $V(0,1,-10)$, determine which points obscure the others when viewed from $V$.
\end{solutionbox}
\begin{answerbox}
$A-V=(5,0,8)$, $B-V=(10,3,19)$, $C-V=(15,-3,7)$. Since $A-V$ has a zero $y$-component, any point collinear with $A$ must also have $P_y-V_y=0$, i.e.\ $P_y=1$ --- but $B_y=4$ and $C_y=-2$, so neither is collinear with $A$. Checking $B$ vs $C$: $(10,3,19)$ vs $(15,-3,7)$ --- opposite-signed $y$-components rule out any positive scalar multiple.
\textbf{Conclusion: no two of $A,B,C$ share a line of sight from $V$ --- none obscures any other, same structure as 2022.}
\end{answerbox}

%=====================================================================
\section{Back-Face / Visible-Surface Definitions}
%=====================================================================
\begin{solutionbox}
\textbf{2024 Q7(a) [2 marks]:} Define visible surface and back-face with necessary examples.
\textbf{Related --- 2024 Q1(c) [3 marks]:} What is hidden surface removal in computer graphics, and what problem does it solve during image generation?
\end{solutionbox}
\begin{answerbox}
\begin{itemize}[noitemsep]
  \item \textbf{Visible surface:} a surface (or part of one) that is actually seen by the viewer from the current viewpoint --- i.e.\ not blocked by any other surface along the line of sight. \emph{Example:} the front face of a cube facing the camera.
  \item \textbf{Back-face:} a polygon whose outward normal points \textbf{away} from the viewer (angle between the normal $\mathbf{N}$ and the view direction $\mathbf{V}$ is $>90°$, i.e.\ $\mathbf{N}\cdot\mathbf{V}<0$ using an outward-pointing convention). \emph{Example:} the far side of a cube facing away from the camera. For a convex, opaque solid, every back-face is automatically hidden, so back-face detection/culling is a cheap first pass before more expensive hidden-surface algorithms (see [[wiki/z-buffer|Z-Buffer, Ch10]]).
  \item \textbf{Hidden-surface removal (related, Q1c):} the process of determining which surfaces (or parts of surfaces) are visible from the viewpoint and discarding/not-rendering the rest. \textbf{Problem it solves:} without it, a rendered image would show every polygon in the scene overlapping incorrectly (e.g.\ surfaces that should be blocked by nearer objects would still appear), producing a visually wrong, non-realistic image — HSR ensures only what would actually be seen ends up on screen.
\end{itemize}
\end{answerbox}

%=====================================================================
\section{Trapezoid-Primitive Polygon Fill (2024 Q7b, Q7c, Q7e)}
%=====================================================================
\begin{notesbox}
No textbook resource for this cluster --- 2024 lecture-only material. The figures in the actual exam (a 6-7 vertex polygon with labeled edges $E_1$--$E_9$ for Q7c, and a separate irregular notched polygon for Q7e) are hand-drawn and only partially legible from the scan. The \textbf{method} below is what earns the marks; apply it directly to your own printed copy of each figure for the exact vertex/edge labels and trapezoid count.
\end{notesbox}

\subsection*{Q7(b) --- Steps of Trapezoid-Primitive Polygon Drawing [2 marks]}
\begin{answerbox}
\begin{enumerate}[noitemsep]
  \item Sort all polygon vertices by $y$-coordinate; identify the distinct $y$-levels.
  \item Draw a horizontal line through every vertex, splitting the polygon into horizontal strips between consecutive distinct $y$-levels.
  \item Each strip is bounded on the left and right by (portions of) polygon edges --- if exactly 2 edges cross a strip, that strip is one \textbf{trapezoid} (a triangle is the degenerate case where one bounding edge shrinks to a point).
  \item If more than 2 edges cross a strip (can happen at a concave notch), split that strip into multiple trapezoids, one per left/right edge-pair.
  \item Fill each trapezoid independently (a trapezoid's left/right boundaries are simple line segments, so each scanline within it needs only 2 boundary-intersection computations) --- this is the appeal of the method: filling reduces to a sequence of simple trapezoid fills instead of one intersection-sorting pass per scanline.
\end{enumerate}
\end{answerbox}

\subsection*{Q7(c) --- Vertex, Edge, and Surface Matrices [2 marks]}
\begin{answerbox}
This is a \textbf{boundary representation}: three linked tables describing a polygon (or, in 3D, a polyhedron) structurally rather than just as a coordinate list.
\begin{itemize}[noitemsep]
  \item \textbf{Vertex table:} one row per vertex, storing its coordinates (e.g.\ $V_1(x_1,y_1),\dots,V_n(x_n,y_n)$).
  \item \textbf{Edge table:} one row per edge, storing the two vertex indices it connects, plus (in 3D) pointers to the one or two surfaces bordering it (e.g.\ $E_1:(V_1,V_2)$, $E_2:(V_2,V_3)$, \dots, following the labels in the given figure).
  \item \textbf{Surface table:} one row per face, listing the ordered sequence of edges (or vertices) bounding it — for a single 2D polygon like the exam figure, one row referencing all of $E_1,\dots,E_9$ in boundary order.
\end{itemize}
\textbf{For the given figure:} read off each vertex position and its connecting edges directly from the diagram (labeled $E_1$--$E_9$ around the boundary, e.g.\ vertex ``$C$'' at the top connects via $E_2$ and $E_3$, vertex ``$F$'' connects via $E_4$ and $E_5$, etc.) and fill the three tables in the format above --- the structure is identical regardless of the specific shape.
\end{answerbox}

\subsection*{Q7(e) --- Number of Trapezoids for the Given Polygon [1 mark]}
\begin{answerbox}
Apply the method from Q7(b): count the number of \textbf{distinct $y$-levels} among the polygon's vertices, draw a horizontal line through each, and count the resulting strips (adding an extra split for any strip crossed by more than 2 edges, which happens at concave notches/zigzags). \textbf{Illustrative example:} a simple convex hexagon with 4 distinct vertex $y$-levels produces $4-1=3$ trapezoids (no notches to force extra splits). For the exam's own (non-convex, zigzag) figure, count the vertex $y$-levels directly off the printed diagram and add one extra trapezoid for each notch that produces more than 2 edge-crossings in its strip.
\end{answerbox}

%=====================================================================
\section{Quick Summary}
%=====================================================================
\begin{center}
\renewcommand{\arraystretch}{1.3}
\begin{tabular}{p{4.6cm}|p{2.2cm}|p{6.8cm}}
\toprule
\textbf{Topic} & \textbf{Years} & \textbf{Method} \\
\midrule
Point-obscures-point & 2020-2023 & $(P-V)$ parallel \& same-sign $\Rightarrow$ same sightline; smaller $t$ obscures larger; 2 of 4 years have \emph{no} occlusion \\
Back-face/visible-surface & 2024 & Back-face: $\mathbf{N}\cdot\mathbf{V}<0$ (points away from viewer); visible surface: unobstructed along sightline \\
Trapezoid-primitive fill & 2024 & Horizontal strips through each vertex $y$-level = trapezoids; extra splits at concave notches \\
Vertex/edge/surface matrices & 2024 & 3 linked tables: vertex coords, edge-to-vertex(+surface) links, surface-to-edge-loop \\
\bottomrule
\end{tabular}
\end{center}

\end{document}
